\section{Discussion and Limitations}
\label{sec:disc}
We report these candidly, per ACM/COPE integrity expectations. (1) \textbf{Storage coupling, built
and remaining.} Leader election is now storage-weighted via on-chain storage-proof records
(\S\ref{sec:consensus},~\ref{sec:audit}): the running consensus is genuinely storage-coupled. Two
related items remain. \emph{Fork-choice} weight is still chain length (we weight \emph{who leads},
not yet \emph{which branch wins}), and missed storage proofs incur election-weight decay rather than
a \emph{bond} slash. Adding storage weight to fork choice and bond slashing for missed proofs are
the priority next steps and would close the gap between the leader-level and branch-level coupling. (2) \textbf{State commitments}: the header
commits transaction and log roots but not a \emph{state root}, so light clients trust state via
replay/snapshot rather than a succinct proof; adding a state (or verkle) root would enable
trust-minimized light clients~\cite{bunz2020flyclient}. (3) \textbf{Data availability} is an
operational reconstruction check, not a sampling proof; DAS~\cite{hallandersen2025foundations}
is complementary. (4) \textbf{Evaluation scale}: a single deployment and workstation, not a
large multi-region benchmark; numbers are feasibility evidence. (5) \textbf{No external security
audit}; the live network is treated as a testnet with no real value. (6) The contract VM is a
Solidity subset. None of these undermines the central result, that history can be fragmented so
no node stores it all, while state stays local and block production is earned by storage, but
each bounds how far the present prototype should be trusted.

\textbf{Applicability.} ShadowLedger is a fit when history grows large relative to active state and
many modest nodes are available to share it: public ledgers, audit logs, and data-preservation
chains, where the archive dwarfs the working set and decentralization is valued over raw
throughput. It is a poor fit when the working state itself is enormous (then keeping state local is
the bottleneck, and state sharding as in OmniLedger~\cite{kokoriskogias2018omniledger} or
RapidChain~\cite{zamani2018rapidchain} is the right tool), or when historical reads are frequent
and latency-critical (the reconstruction cost of \S\ref{sec:eval} would dominate). The design also
assumes enough independent nodes that the $r$ holders of a shard fail independently; in a tiny or
collusive network that assumption, and with it the availability argument of \S\ref{sec:threat},
weakens. We therefore present ShadowLedger as a point in the design space, not a universal
replacement for replicated ledgers.

\textbf{On head-to-head benchmarking.} We deliberately do not report wall-clock comparisons against
deployed systems such as Filecoin, Chia, or other coded-blockchain prototypes. Such numbers would
be measured on different hardware, networks, and workloads, so a cross-system performance table
would be apples-to-oranges and, by our own integrity standard, misleading; reproducing each system
in a controlled common testbed is a substantial effort beyond this paper. Instead we compare where
comparison is fair: analytically and under identical assumptions against full replication and naive
distributed replication (\S\ref{sec:eval}), and qualitatively against prior coded-storage designs
(\S\ref{sec:related}). The distinction worth stating plainly is not a larger constant-factor saving
than prior coded nodes, which report substantial per-node reductions at a \emph{fixed} network size;
it is that ShadowLedger's reduction \emph{scales} as $O(1/n)$ with the membership and is produced by
the same mechanism that grants block-production rights, rather than by a storage layer bolted onto an
unchanged consensus. A controlled multi-system benchmark on shared infrastructure is the natural next
empirical step.
